% ============================================================================
%  Chapitre 2, partie II : Repères et rotations
%  (manuscrit P1 p. 20 à 24,
%  P2 p. 1 à 4)
% ============================================================================

% ----------------------------------------------------------------------------
\section{Repères et rotations}
% ----------------------------------------------------------------------------

On reprend la convention des repères : $\rep{s}$ lié au sol et
$\rep{b}$ lié au solide, orthonormés directs. Pour un vecteur
$\vect{u}$, on note $\vect{u}_s$ et $\vect{u}_b$ ses colonnes de
composantes dans ces deux repères, et
\[
    \vect{u}_s = \mat{R}\,\vect{u}_b,
    \qquad \mat{R} := \mat{R}_{b\to s}
\]

% ----------------------------------------------------------------------------
\subsection{Représentation d'une orientation : \texorpdfstring{$\SO$}{SO(3)}
et \texorpdfstring{$S^3$}{S3}}

Une rotation de l'espace peut être représentée par une matrice de
\[
    \SO = \{\mat{R}\in\mathcal{M}_3(\R) :
    \mat{R}^T\mat{R} = \Id,\ \det\mat{R} = 1\}
\]
\begin{itemize}
    \item $\mat{R}^T\mat{R} = \Id \longleftrightarrow
        \mat{R}\in\mathrm{O}(3)$ : la rotation conserve le produit
        scalaire ;
    \item $\det\mat{R} = 1$ : elle conserve le volume et l'orientation.
\end{itemize}

Une autre représentation utilise les \emph{quaternions unitaires}
\[
    S^3 = \{\quat{q} = (q_0,q_1,q_2,q_3)\in\R^4 : \|\quat{q}\|^2 = 1\}
\]
$S^3$ désigne la sphère de dimension trois plongée dans $\R^4$. Cette
représentation n'est pas unique : $\mat{R}(\quat{q})
= \mat{R}(-\quat{q})$, donc chaque rotation correspond à deux
quaternions opposés (revêtement double de $\SO$ par $S^3$).

% ----------------------------------------------------------------------------
\subsection{Matrices de rotation et cinématique sur \texorpdfstring{$\SO$}{SO(3)}}

\subsubsection{Matrice de rotation}

\begin{definition}[Matrice de rotation]
    Soient $\vect{b}_1,\vect{b}_2,\vect{b}_3$ les axes du solide
    exprimés dans le repère sol. La matrice de rotation est
    \[
        \mat{R} = \begin{pmatrix}
            \vect{b}_1 & \vect{b}_2 & \vect{b}_3
        \end{pmatrix},
        \qquad
        \mat{R}^T\mat{R} = \Id,
        \qquad \det\mat{R} = 1
    \]
    Ainsi $\mat{R}^{-1} = \mat{R}^T$ et $\vect{u}_b
    = \mat{R}^T\vect{u}_s$.
\end{definition}

\subsubsection{Matrices antisymétriques et vitesse angulaire}

On note $\mathcal{A}_3(\R) = \{\mat{A}\in\mathcal{M}_3(\R) :
\mat{A}^T = -\mat{A}\}$ l'ensemble des matrices antisymétriques. À
$\vect{u}\in\R^3$ on associe
\[
    \hatx{\vect{u}} = \begin{pmatrix}
        0 & -u_3 & u_2\\
        u_3 & 0 & -u_1\\
        -u_2 & u_1 & 0
    \end{pmatrix},
    \qquad
    \hatx{\vect{u}}\vect{v} = \vect{u}\wedge\vect{v}
\]

\begin{proposition}[Cinématique d'une matrice de rotation]
    Si $t\mapsto\mat{R}(t)$ est une courbe dérivable dans $\SO$, il
    existe un unique $\vect{\omega}(t)\in\R^3$ tel que
    \[
        \dot{\mat{R}} = \mat{R}\,\hatx{\vect{\omega}}
    \]
\end{proposition}

\begin{demonstration}
    En dérivant $\mat{R}^T\mat{R} = \Id$,
    \[
        \dot{\mat{R}}^T\mat{R} + \mat{R}^T\dot{\mat{R}} = \mat{0}
        \longleftrightarrow
        \bigl(\mat{R}^T\dot{\mat{R}}\bigr)^T
        = -\mat{R}^T\dot{\mat{R}}
    \]
    donc $\mat{R}^T\dot{\mat{R}}\in\mathcal{A}_3(\R)$, et il existe un
    unique $\vect{\omega}$ tel que $\mat{R}^T\dot{\mat{R}}
    = \hatx{\vect{\omega}}$
\end{demonstration}

$\vect{\omega}$ est la vitesse angulaire exprimée dans le corps, et
$\vect{\Omega} = \mat{R}\,\vect{\omega}$ la même vitesse exprimée dans
le sol. Comme $\mat{R}(\vect{u}\wedge\vect{v})
= (\mat{R}\vect{u})\wedge(\mat{R}\vect{v})$, on a
$\mat{R}\,\hatx{\vect{u}} = \hatx{\mat{R}\vect{u}}\,\mat{R}$, d'où
\[
    \dot{\mat{R}} = \mat{R}\,\hatx{\mat{R}^T\vect{\Omega}}
    = \mat{R}\mat{R}^T\hatx{\vect{\Omega}}\,\mat{R}
    = \hatx{\vect{\Omega}}\,\mat{R}
\]

\subsubsection{Formule de transport}

\begin{theoreme}[Formule de transport]
    Si $\vect{u}_s(t) = \mat{R}(t)\,\vect{u}_b(t)$, alors
    \[
        \dot{\vect{u}}_s
        = \mat{R}\bigl(\dot{\vect{u}}_b + \vect{\omega}\wedge\vect{u}_b\bigr)
    \]
\end{theoreme}

\begin{preuve}
    $\dot{\vect{u}}_s = \dot{\mat{R}}\,\vect{u}_b
    + \mat{R}\,\dot{\vect{u}}_b
    = \mat{R}\bigl(\dot{\vect{u}}_b
    + \hatx{\vect{\omega}}\vect{u}_b\bigr)
    = \mat{R}\bigl(\dot{\vect{u}}_b + \vect{\omega}\wedge\vect{u}_b\bigr)$
\end{preuve}

\subsubsection{Équations d'Euler}

Dans le repère corps, $\vect{L}_b = \mat{I}_b\,\vect{\omega}$ avec
$\mat{I}_b$ constante. La formule de transport appliquée à
$\vect{L}_s = \mat{R}\,\vect{L}_b$ donne
\[
    \mat{R}^T\dot{\vect{L}}_s
    = \mat{I}_b\,\dot{\vect{\omega}}
    + \vect{\omega}\wedge(\mat{I}_b\,\vect{\omega})
\]

\begin{theoreme}[Équations d'Euler]
    Avec le théorème du moment cinétique, en notant
    $\vect{M}_b^{\mathrm{ext}}$ le moment extérieur exprimé dans le
    corps,
    \[
        \mat{I}_b\,\dot{\vect{\omega}}
        + \underbrace{\vect{\omega}\wedge(\mat{I}_b\,\vect{\omega})}
        _{\text{terme gyroscopique}}
        = \vect{M}_b^{\mathrm{ext}}
    \]
\end{theoreme}

Le terme gyroscopique apparaît dès que le solide tourne par rapport au
sol. Pour un solide anisotrope, $\vect{\omega}$ n'est pas forcément parallèle à
$\vect{L}_b$, et $\vect{\omega}\wedge(\mat{I}_b\vect{\omega})$ décrit
la contribution de la rotation du solide au changement de direction du
moment cinétique :
\[
    \mat{I}_b\,\dot{\vect{\omega}}
    = \vect{M}_b^{\mathrm{ext}}
    - \vect{\omega}\wedge(\mat{I}_b\,\vect{\omega})
\]

\begin{remarque}[Interprétation gyroscopique]
    Le terme $\vect{\omega}\wedge\vect{L}_b$ ne conserve rien par
    lui-même, il traduit seulement la rotation de $\rep{b}$. En
    l'absence de moment extérieur, c'est $\vect{L}_s = \mat{R}\,
    \vect{L}_b$ qui est constant, alors que ses composantes
    $\vect{L}_b$ varient
    \[
        \dot{\vect{L}}_b = -\,\vect{\omega}\wedge\vect{L}_b
    \]
    \begin{itemize}
        \item Si $\vect{\omega}$ est porté par un axe principal,
            $\vect{\omega}\parallel\mat{I}_b\vect{\omega}$, le terme
            s'annule et la rotation se maintient : c'est la toupie qui
            continue de tourner.
        \item Pour une toupie pesante, le poids exerce un moment
            perpendiculaire à $\vect{L}$, avec une grande rotation
            propre, $\dot{\vect{L}} = \vect{M}$ fait tourner l'axe
            autour de la verticale (précession) au lieu de le faire
            basculer.
    \end{itemize}
\end{remarque}

% ----------------------------------------------------------------------------
\subsection{Quaternions unitaires et \texorpdfstring{$S^3$}{S3}}

\begin{definition}[Quaternion]
    Un quaternion est $\quat{q} = (q_0,\vect{q})\in\R\times\R^3$, de
    partie scalaire $q_0$ et de partie vectorielle $\vect{q}$. Le
    produit de Hamilton est
    \[
        \quat{q}\quatmul\quat{p}
        = \bigl(q_0p_0 - \vect{q}^T\vect{p},\
        q_0\vect{p} + p_0\vect{q} + \vect{q}\wedge\vect{p}\bigr)
    \]
\end{definition}

Ce produit n'est pas commutatif, à cause du terme
$\vect{q}\wedge\vect{p}$. On définit le conjugué et la norme
\[
    \quat{q}^* = (q_0,-\vect{q}),
    \qquad
    \|\quat{q}\|^2 = q_0^2 + \|\vect{q}\|^2,
    \qquad
    \quat{q}\quatmul\quat{q}^* = (\|\quat{q}\|^2,\vect{0}),
    \qquad
    \quat{q}^{-1} = \frac{\quat{q}^*}{\|\quat{q}\|^2}
\]
et la norme est multiplicative : $\|\quat{q}\quatmul\quat{p}\|
= \|\quat{q}\|\,\|\quat{p}\|$

\subsubsection{Représentation d'une rotation}

À un vecteur $\vect{u}\in\R^3$ on associe le quaternion pur
$(0,\vect{u})$. Pour un quaternion unitaire $\quat{q}$, on définit
\[
    (0,\vect{u}_s) = \quat{q}\quatmul(0,\vect{u}_b)\quatmul\quat{q}^*
\]

\begin{proposition}[Matrice de rotation d'un quaternion]
    \[
        \vect{u}_s = (q_0^2 - \|\vect{q}\|^2)\,\vect{u}_b
        + 2\,\vect{q}\,(\vect{q}^T\vect{u}_b)
        + 2\,q_0\,(\vect{q}\wedge\vect{u}_b)
    \]
    La matrice associée est donc
    \[
        \mat{R}(\quat{q}) = (q_0^2 - \|\vect{q}\|^2)\,\Id
        + 2\,\vect{q}\vect{q}^T + 2\,q_0\,\hatx{\vect{q}}
    \]
\end{proposition}

Pour une rotation d'angle $\theta$ autour de l'axe unitaire $\vect{u}$,
\[
    \quat{q} = \Bigl(\cos\frac{\theta}{2},\
    \vect{u}\,\sin\frac{\theta}{2}\Bigr)
\]
Comme $\mat{R}(\quat{q}) = \mat{R}(-\quat{q})$ et
$\mat{R}(\quat{q}\quatmul\quat{p})
= \mat{R}(\quat{q})\,\mat{R}(\quat{p})$, l'application
$\quat{q}\mapsto\mat{R}(\quat{q})$ associe à chaque rotation deux
quaternions opposés.

\subsubsection{Cinématique quaternionique}

On suppose que $\quat{q}(t)$ représente l'orientation du corps vers le
sol et que $\vect{\omega}$ est la vitesse angulaire exprimée dans le
corps. Pendant un petit intervalle $h$, le quaternion de rotation
relative est
\[
    \delta\quat{q} = \Bigl(1,\ \frac{h}{2}\,\vect{\omega}\Bigr) + o(h)
\]
et $\quat{q}(t+h) = \quat{q}(t)\quatmul\delta\quat{q}$, d'où
\[
    \dot{\quat{q}}(t)
    = \lim_{h\to 0}\frac{\quat{q}(t+h) - \quat{q}(t)}{h}
    = \tfrac12\,\quat{q}\quatmul(0,\vect{\omega})
\]

\begin{proposition}[Cinématique quaternionique]
    \[
        \dot{\quat{q}} = \tfrac12\,\quat{q}\quatmul(0,\vect{\omega})
        \qquad\longleftrightarrow\qquad
        \left\{
        \begin{aligned}
            \dot q_0 &= -\tfrac12\,\vect{q}^T\vect{\omega}\\
            \dot{\vect{q}} &= \tfrac12\bigl(q_0\,\vect{\omega}
            + \vect{q}\wedge\vect{\omega}\bigr)
        \end{aligned}
        \right.
    \]
    Avec la vitesse angulaire exprimée dans le sol,
    $\vect{\Omega} = \mat{R}(\quat{q})\,\vect{\omega}$, la même
    équation s'écrit
    \[
        \dot{\quat{q}} = \tfrac12\,(0,\vect{\Omega})\quatmul\quat{q}
    \]
\end{proposition}

\begin{proposition}[Conservation de la norme]
    \[
        \ddt\|\quat{q}\|^2
        = 2\,q_0\dot q_0 + 2\,\vect{q}^T\dot{\vect{q}}
        = -q_0\,\vect{q}^T\vect{\omega} + q_0\,\vect{q}^T\vect{\omega}
        + \vect{q}^T(\vect{q}\wedge\vect{\omega}) = 0
    \]
    Un quaternion unitaire conserve sa norme au cours de son mouvement.
\end{proposition}
